\scanpage{9}
($\operatorname{Ob}\mathcal X=X,\ \operatorname{Fl}\mathcal X=G\times X\ldots$), \deleted{\unsure{on prend}}\inserted{$\alpha_X$ l'}application canonique
\[
K_1=\operatorname{Fl}\mathcal X=\unsure{G|C}\times X\xrightarrow{\mathrm{pr}_1}G
\]
On pose
\[
\widetilde{\unsure{\mathscr P|P}}(X)=\operatorname{Rep}(\mathcal X;\mathscr P,\alpha_X)=\operatorname{Rep}(X,\mathscr P,\alpha_{\unsure{x|X}})
\]
avec la contravariance p.r. à $X$ déjà vue. On trouve que c'est un 2-champ, qui est $=$ une 2-gerbe liée par $M_{(B_G)}$ \deleted{\unsure{(muni de|construit}} (défini par l'op. de $G$ sur $M$).

\deleted{\unsure{3.}} \uline{\unsure{Variance|Varions|Variante}} $\widetilde{\mathscr P}$ est covariant en $\mathscr P$ dans un sens évident \ldots\ \deleted{\unsure{On a un fon}}\inserted{En fait \unsure{un|-}} un morphisme de la 3-catégorie des \deleted{3. Fixons} $\mathscr P$ dans la 3-catégorie des 2-gerbes sur des $B_G$ variables.

3. On a ainsi une équivalence de 3-catégories. Plus \unsure{spécifiquement|explicitement}, $G$ et le $G$-module $M$ étant fixés, on a une équivalence entre la 3-catégorie des Gr-catégories \deleted{de Picard} épinglées par $(G,M)$, et la 3-catégorie des 2-gerbes sur $B_G$ liées par $M_{(B_G)}$

Soient
\[
\begin{aligned}
&\mathscr P &&\text{Gr-catégorie }\deleted{\unsure{\text{définie}|\text{d'invariants}}}\inserted{\text{épinglée par}} &&G,M\\
&\mathscr P' &&\rule{5em}{.3pt} &&G',M'
\end{aligned}
\]
\[
G\overset{u}{\unsure{\relbar|\longrightarrow}}G',\quad M\overset{\unsure{v|o}}{\relbar}M'\qquad\text{hom compatibles aux opérations}
\]
Alors les $\otimes$ hom $\mathscr P\relbar\mathscr P'$ compatibles avec $G\relbar G', M\relbar M'$ forment une \uline{2-catégorie}
% archive sheet number: 8
